\begin{trapbox}
	\begin{description}[style=nextline, leftmargin=2.8em, labelsep=0.5em, itemsep=0.35em]
		\item[\textbf{\textcolor{CTrap}{易错点一（SAS夹角条件）}}] 认为两边成比例且任意一角相等就可判定相似，忽略夹角要求。
		\item[\textbf{\textcolor{CTrap}{正确理解（需夹角相等）：}}] 在$\triangle ABC$与$\triangle DEF$中，若$\frac{AB}{DE}=\frac{BC}{EF}$且$\angle B=\angle E$（$\angle B$是$AB$与$BC$的夹角），才能用SAS判定；若$\frac{AB}{DE}=\frac{BC}{EF}$但$\angle A=\angle D$（非夹角），不能直接判定。
		\item[\textbf{\textcolor{CTrap}{错因分析（非夹角反例）：}}] 夹角不相等时，边比例固定但夹角可变，三角形形状不唯一，反例容易构造。
	\end{description}

	\medskip
	\begin{description}[style=nextline, leftmargin=2.8em, labelsep=0.5em, itemsep=0.35em]
		\item[\textbf{\textcolor{CTrap}{易错点二（HL直角限制）}}] 在非直角三角形中滥用HL判定，或者忘记需要直角条件。
		\item[\textbf{\textcolor{CTrap}{正确理解（直角前提）：}}] HL判定仅适用于直角三角形，必须声明$\angle C=\angle C'=90^\circ$，然后斜边和一条直角边成比例。
		\item[\textbf{\textcolor{CTrap}{错因分析（无直角行不通）：}}] HL本质是通过勾股定理转化为SSS；无直角则无法推出第三边比例，判定不成立。
	\end{description}

	\medskip
	\begin{description}[style=nextline, leftmargin=2.8em, labelsep=0.5em, itemsep=0.35em]
		\item[\textbf{\textcolor{CTrap}{易错点三（对应边对应）}}] 在利用相似列比例式时，对应边找错，导致方程错误。
		\item[\textbf{\textcolor{CTrap}{正确理解（顶点顺序）：}}] 要严格按相似记作“$\triangle ABC\sim\triangle DEF$”，则对应边为$AB$与$DE$，$BC$与$EF$，$CA$与$FD$；注意字母顺序。
		\item[\textbf{\textcolor{CTrap}{错因分析（混乱顺序）：}}] 学生常将$AB$对应$EF$等，不作严格对应，导致比例关系乱，求解必然错误。
	\end{description}
\end{trapbox}
